\chapter{Related Work}\label{ch:related-work}

Repository-level program comprehension requires locating relevant source across multiple files and reconstructing how the responsible implementation elements interact to produce observable behaviour. Research has addressed parts of this problem through approaches with different aims. Code localisation seeks relevant files or functions (\textcite{Zhou2012BugLocator}), whereas repository question answering retrieves material to support an answer (\textcite{Peng2026SWEQA}). Structural navigation uses relationships between code elements to guide inspection (\textcite{Petrenko2013DepIR}), while language-model-directed exploration delegates inspection choices to a model (\textcite{Kang2024AutoFL}). Other work addresses the traceability of how material is obtained through provenance (\textcite{W3C2013PROV}), or evaluates whether cited material supports generated claims (\textcite{Gao2023ALCE}).

Their evaluations consequently ask different questions: whether relevant files or functions rank highly, whether navigation reaches the responsible code, or whether an answer's claims are supported by its cited passages. This chapter compares these approaches and identifies what remains insufficiently examined about constructing and evaluating repository evidence.

\section{From Code Localisation to Repository-Level Evidence Retrieval}\label{sec:from-code-localisation-to-repository-level-evidence-retrieval}

Empirical studies of software development establish that repository understanding involves several related information needs. \textcite{Ko2007InformationNeeds} find that developers ask where code is located, what it does, why it behaves in a particular way, and how its parts interact. \textcite{Sillito2008ProgrammingQuestions} similarly describe a progression from locating an initial focus to understanding larger structures, dependencies, and implications during a change task. These findings distinguish the identification of relevant source from the broader task of explaining how that source contributes to system behaviour.

Finding the source that corresponds to these information needs has therefore been a central problem in program-comprehension research. Two central applications of this source-identification task are concept localisation and bug localisation. \textcite{Marcus2004ConceptLocation} represent source documents in a latent semantic space, allowing natural-language concepts to be associated with code despite differences in vocabulary. BugLocator ranks source files using text similarity to a bug report together with evidence from previously resolved reports (\textcite{Zhou2012BugLocator}). BLUiR instead preserves the structure of bug reports and source files by retrieving across their constituent fields rather than treating each artifact as undifferentiated text (\textcite{Saha2013BLUiR}). These approaches demonstrate how semantic representation, development history, and document structure can improve the ranking of sources associated with a concept or reported problem.

The target of this research nevertheless remains localisation. Success is commonly determined by whether a known relevant file or element appears sufficiently high in the ranking. This is necessary for repository understanding, but it is not sufficient when an explanation depends on several snippets. A correctly ranked file may not identify the responsible function, while an isolated function may not reveal its caller, state transition, or downstream effect. The distinction is therefore not simply between successful and unsuccessful localisation, but between retrieving an individually relevant result and assembling the multiple snippets needed to explain how a mechanism operates.

Repository-level question answering expands the scope beyond an individual source location. SWE-QA (\textcite{Peng2026SWEQA}) indexes functions, classes, and their dependencies, then combines semantic similarity with structural dependency analysis to retrieve relevant code, documentation, and architectural metadata. Text similarity helps identify code relevant to the question, while dependency information helps extend the retrieved context across the source relationships required for cross-file and multi-hop questions. Because an entire repository cannot receive equal attention within a bounded model context, this broader scope also creates a selection problem: retrieval must prioritise which snippets to include. Additionally, repository-question benchmarks measure retrieval and answer quality, but aggregate scores do not show whether the resulting context contains all snippets required to explain a multi-file mechanism. This limitation motivates evaluating repository retrieval as the construction of a complete evidence set, rather than solely as the ranking of individual results.

\section{Structural and Graph-Guided Navigation}\label{sec:structural-and-graph-guided-navigation}

When several snippets are required, textual relevance alone provides limited guidance about how they are connected. Earlier code-search systems therefore combined textual relevance with structural context (\textcite{McMillan2011Portfolio,Petrenko2013DepIR}). Portfolio (\textcite{McMillan2011Portfolio}) retrieves functions and presents their usages, using call-graph relationships to place a matching function within its surrounding code. DepIR (\textcite{Petrenko2013DepIR}) combines lexical search over code identifiers and comments with navigation over statically inferred program dependencies. Across 50 historical changes in five systems, its evaluation reports lower localisation effort than either text retrieval or dependency navigation alone.

Recent language-model systems make repository structure available during model-directed navigation. CodexGraph stores code entities and relationships in a graph database and allows an LLM to construct graph queries for structure-aware repository access (\textcite{Liu2025CodexGraph}). LocAgent builds a heterogeneous graph containing files, classes, functions, imports, invocations, and inheritance, then uses a language-model agent to traverse it for multi-hop code localisation (\textcite{Chen2025LocAgent}). Compared with flat chunk retrieval, these representations expose named code entities and explicit routes between them.

Structural reachability is not equivalent to semantic evidence. A call, reference, import, or containment edge establishes a relation inferred by an analysis, but it does not establish that the connected node answers the user's question. A general-purpose utility function may be called from many parts of a repository without helping explain the particular behaviour in question. Static analysis can also omit runtime registration, dynamic dispatch, reflection, generated artifacts, or relationships that a language-specific analyser does not represent. Structural navigation can consequently recover candidates missed by textual ranking, but the reached source still requires inspection and qualification.

The literature (\textcite{McMillan2011Portfolio,Petrenko2013DepIR}) therefore suggests two distinct contributions of structure. It can contextualise a textual match within its enclosing code entity, and it can expand retrieval towards related source that was not highly ranked by the original query. Evaluating only whether a graph-assisted system eventually finds a target location does not distinguish these contributions or show whether graph-discovered material survives into the final context. This leads to the question of who decides what to inspect next and when exploration should stop.

\section{Agentic Repository Exploration}\label{sec:agentic-repository-exploration}

Exploration directed by a language model makes repository retrieval adaptive. ReAct establishes the general pattern of interleaving model reasoning with actions that produce new observations (\textcite{Yao2023ReAct}). Applied to repositories, such actions can include searching for terms, opening source ranges, locating symbols, following references, reformulating a query, and deciding whether the current material is sufficient. The exploration path can respond to newly observed identifiers and relationships rather than being fixed before retrieval begins.

AutoFL illustrates this allocation in fault localisation. It allows an LLM to inspect repository information through tools, select likely fault locations, and produce an explanation of its conclusion (\textcite{Kang2024AutoFL}). LocAgent similarly lets the model choose graph-navigation steps according to the entities already observed (\textcite{Chen2025LocAgent}). Although their immediate evaluation target is localisation, both are relevant to repository evidence acquisition because the model determines which observations lead to further inspection.

This flexibility creates an attribution problem. A useful final source set may result from the initial query, subsequent tool use, structural navigation, or model inference. Retaining only the final locations therefore obscures how the evidence was obtained and why potentially relevant source was omitted. This limitation connects adaptive repository exploration to the broader concerns of grounding and provenance.

\section{Grounding and Provenance}\label{sec:grounding-and-provenance}

The need to account for how evidence was obtained connects adaptive exploration to grounding. Retrieval-augmented language modelling demonstrates how external documents can condition an answer instead of leaving it entirely to model parameters (\textcite{Lewis2020RAG}). Providing context does not guarantee that every generated claim follows from that context. ALCE evaluates answers along citation correctness and completeness and shows that plausible responses can still contain claims without adequate citation support (\textcite{Gao2023ALCE}). For repository understanding, attaching a file reference to an explanation is correspondingly weaker than showing that a particular source range supports the associated implementation claim.

Provenance provides a way to record how supporting material was obtained. The W3C PROV model describes entities, activities, agents, and the relationships involved in producing an artefact (\textcite{W3C2013PROV}). In repository retrieval, analogous records can associate a source range with the query and retrieval channel that found it, the structural operation that resolved or expanded it, and the decisions that retained or rejected it. Provenance does not make an incorrect judgement correct, but it allows that judgement and its inputs to be inspected.

Existing approaches expose different fragments of this history. Search systems expose rankings, graph systems expose paths or neighbours, agents expose tool trajectories, and citation systems expose final claim-to-source links. These records answer different questions. A citation can show which source is offered in support of a final statement without showing how that source entered the context. An agent trajectory can show that a file was opened without showing whether its contents were later admitted as evidence. A ranked list can show initial relevance without showing later survival.

The underexplored boundary is therefore the lifecycle between retrieval and final support. Evaluating that boundary requires distinguishing source that was never retrieved from source that was later transformed, deferred, rejected, recovered, or selected. It also requires enough provenance to reconstruct which model and application decisions affected each candidate and how bounded resources constrained those decisions. This lifecycle perspective connects repository retrieval with source-grounded assistance: the final explanation is only as inspectable as the evidence process that produced its supporting context.

\section{Learning-Oriented Assistance}\label{sec:learning-oriented-assistance}

Repository assistance can present the source directly or use it to support active engagement with an implementation. Research on self-explanation provides an educational basis for the latter approach. The meta-analysis by \textcite{Bisra2018SelfExplanation} reports an overall positive, though heterogeneous, effect from prompting learners to explain material to themselves. In source-code comprehension, \textcite{Tamang2021Socratic} compare free self-explanations with Socratic tutoring explanations, directly examining prompts that ask learners to articulate relationships in code.

These findings motivate guided questions, staged hints, and feedback that preserve part of the reasoning task for the learner. They also make evidence quality consequential because unsupported guidance can promote an incorrect account of program behaviour. Source grounding can constrain what an assistant asks the learner to infer, but it cannot establish that learning occurs. The educational literature therefore motivates the interaction design, while claims about learning outcomes would require a separate user study. The empirical subject of this thesis remains repository retrieval and evidence construction.

\section{Comparative Synthesis and Research Gap}\label{sec:comparative-synthesis-and-research-gap}

The reviewed work progresses from locating source to assembling support for an explanation, but each line makes a different part of that progression visible.

\begin{table}[htbp]
\centering
\caption{Comparative synthesis and research gap}
\label{tab:comparative-synthesis-and-research-gap-1}

\begingroup
\small
\setlength{\tabcolsep}{4pt}
\renewcommand{\arraystretch}{1.15}

\begin{tabularx}{\linewidth}{
  @{}L{0.20\linewidth}YY@{}
}
\toprule
\textbf{Research line}
& \textbf{Evidence unit and retrieval signal}
& \textbf{Decision ownership, evaluation, and visibility} \\
\midrule

Concept and bug localisation
& \textit{Unit:} File or source document.
  \par\textit{Signal:} Lexical, semantic, historical, and field-specific information.
& \textit{Ownership:} Fixed retrieval pipeline.
  \par\textit{Evaluation:} Known-location ranking.
  \par\textit{Visibility:} Initial ranking. \\

Repository-level retrieval and QA
& \textit{Unit:} Repository or cross-file context.
  \par\textit{Signal:} Dense retrieval and context augmentation.
& \textit{Ownership:} Pipeline or model policy.
  \par\textit{Evaluation:} Retrieval and answer quality.
  \par\textit{Visibility:} Retrieved and final context. \\

Structural navigation
& \textit{Unit:} Function, owner, or graph node.
  \par\textit{Signal:} Calls, references, dependencies, and containment.
& \textit{Ownership:} User, pipeline, or graph-enabled agent.
  \par\textit{Evaluation:} Location recovery and navigation effort.
  \par\textit{Visibility:} Graph paths and final locations. \\

Agentic repository exploration
& \textit{Unit:} Source ranges and tool observations.
  \par\textit{Signal:} Iterative search, inspection, and navigation.
& \textit{Ownership:} Predominantly model-directed.
  \par\textit{Evaluation:} Localisation, explanation quality, and cost.
  \par\textit{Visibility:} Action trajectory. \\

Grounded assistance
& \textit{Unit:} Source passage and explanation claim.
  \par\textit{Signal:} Retrieval plus source attribution.
& \textit{Ownership:} Mixed.
  \par\textit{Evaluation:} Answer and citation support.
  \par\textit{Visibility:} Final claim-to-source link. \\

\bottomrule
\end{tabularx}

\endgroup
\end{table}

Prior work establishes that semantic matching can bridge vocabulary differences, repository context can support broader questions, structural relationships can recover poorly ranked source, and model-directed exploration can adapt to discoveries made during search. Provenance and citation research further show the importance of making support inspectable. These contributions are complementary, but their evaluations usually emphasise an initial ranking, final location, action outcome, or cited answer.

The remaining gap concerns the controlled construction of mechanism evidence between those endpoints. There is limited evaluation of a pipeline that combines hybrid retrieval, structural-owner resolution, semantic qualification, and bounded iterative navigation while exposing how each evidence item is admitted, transformed, deferred, rejected, recovered, or selected. This matters because a final miss does not reveal whether the responsible source was absent from retrieval or lost later, and a final hit does not establish that the selected evidence explains the required cross-file relationships.

\begin{samepage}
This thesis addresses that gap by treating repository evidence as a lifecycle with explicit provenance and bounded decision ownership. It evaluates terminal localisation, the survival of exact required source, and whether all required snippets remain available together, including cases in which those snippets form a connected mechanism. It also compares an application-controlled pipeline with agentic repository exploration in terms of evidence completeness, stability, and cost. Its contribution is not a new claim that text search, graphs, or agents can locate code independently. It is an investigation of how those capabilities can be combined and evaluated as an auditable evidence-construction process for unfamiliar-codebase understanding.
\end{samepage}
